\begin{problem}[第37届全国中学生物理竞赛决赛][散裂中子源与中子慢化]{125}
    我国大科学装置散裂中子源于2018年建成并投入使用，它在诸多领域有广泛的应用。历史上，查德威克在1932年首次确认了中子的存在并测出了它的质量；哈恩等人在1939年发现用中子轰击铀原子核可使其分裂，同时放出中子，引发链式反应。为了使链式反应能够持续可控地进行，可通过弹性碰撞使铀核放出的中子慢化。不考虑相对论效应。
    \begin{subproblem}
        查德威克用中子（质量为$m$）轰击质量为$m _ { 1 }$的静止靶核（氢核$\mathrm{H}$或氮核$^{14}\mathrm{N}$，质量为$m _ { \mathrm { H } }$或$14 m _ { \mathrm { H } }$），观察它们的运动。设靶核的出射动量与入射中子的初动量之间的夹角为$\alpha$，试导出此时靶核的出射速率$v _ { 1 }$和中子的末速率$v$分别与中子初速率$v _ { 0 }$之间的关系。该实验测得氢核的最大出射速率为$3.30 \times 10^{7} \mathrm{m/s}$，氮核的最大出射速率为$4.70 \times 10^{6} \mathrm{m/s}$，求$m$和$v _ { 0 }$的值。
    \end{subproblem}
    \begin{subproblem}
        在上述实验中一个氮核也可能受到一束中子的连续撞击。假设氮核开始时是静止的，每次与之相碰的中子的速率都是$v _ { 0 }$，碰撞都使得氮核速率的增量最大。试计算经过多少次碰撞后氮核的动能与中子的初动能近似相等？
    \end{subproblem}
    \begin{subproblem}
        设经过多次碰撞被减速的中子处于热平衡状态，其速率满足麦克斯韦分布
        \begin{equation}
            f(v) = 4 \pi \left(\frac {m}{2 \pi k _ {\mathrm{B}} T}\right) ^ {3 / 2} v ^ {2} \mathrm{e} ^ {- \frac {m v ^ {2}}{2 k _ {\mathrm{B}} T}},
            \label{eq:1.p125}
        \end{equation}
        这里$k _ { \mathrm { B } }$为玻尔兹曼常量，$T$为系统的绝对温度。试计算在热平衡时，中子的最概然速率所对应的动能和最概然动能。
    \end{subproblem}
\end{problem}